% The national range of the name Adjap: 30 blocks laid end to end from 1000 (Eq. 6), each with the numbers issued at
% first issue, the rest of its main part and its tail reserve (Eq. 9), then the spare. Data: adjap-blocks.dat, written by tools/paper-figures.py.
\documentclass[tikz,border=3pt]{standalone}
\usepackage{fontspec}
\setmainfont{LibertinusSans}[Extension=.otf,UprightFont=*-Regular,BoldFont=*-Bold,ItalicFont=*-Italic]
\usepackage{unicode-math}
\setmathfont{LibertinusMath-Regular.otf}
\usepackage{pgfplots}
\pgfplotsset{compat=1.18}
\usepackage{pgfplotstable}
% English by default; the French paper compiles <name>-fr.tex, which sets \fr to 1 and inputs this file.
\providecommand{\fr}{0}
\newcommand{\T}[2]{\ifnum\fr=1 #2\else #1\fi}
\ifnum\fr=1 \pgfkeys{/pgf/number format/.cd, use comma, 1000 sep={\,}}\else \pgfkeys{/pgf/number format/.cd, 1000 sep={,}}\fi
\definecolor{ink}{HTML}{00664F}\definecolor{flaggreen}{HTML}{007A5E}\definecolor{flagred}{HTML}{CE1126}
\definecolor{flagyellow}{HTML}{FCD116}\definecolor{amber}{HTML}{C98A00}
\usetikzlibrary{patterns}
\begin{document}
\begin{tikzpicture}[x=1cm,y=1cm]
\pgfplotstableread{adjap-blocks.dat}\blocks
\def\W{14}
\pgfmathsetmacro{\sx}{\W/9000}
% the whole range 1000-9999
\fill[gray!12] (0,2.8) rectangle (\W,3.4);
\fill[flaggreen!70] (0,2.8) rectangle ({(4950-1000)*\sx},3.4);
\fill[flagred!25] ({(9100-1000)*\sx},2.8) rectangle (\W,3.4);
\draw[gray] (0,2.8) rectangle (\W,3.4);
\node[font=\footnotesize, anchor=south] at ({(4950-1000)*\sx/2},3.4) {\T{30 blocks, 1000–4949}{30 blocs, 1000–4949}};
\node[font=\footnotesize, anchor=south] at ({(7025-1000)*\sx},3.4) {\T{spare: new areas and extensions}{réserve : nouvelles zones et extensions}};
\node[font=\footnotesize, anchor=south, text=flagred] at ({(9550-1000)*\sx},3.4) {\T{top tenth}{dixième supérieur}};
\foreach \t in {1000,2000,...,9000} {\draw[gray] ({(\t-1000)*\sx},2.8) -- ({(\t-1000)*\sx},2.7) node[below, font=\scriptsize] {\t};}
\draw[gray] (\W,2.8) -- (\W,2.7) node[below, font=\scriptsize] {9999};
% zoom on 1000-4949
\pgfmathsetmacro{\zx}{\W/3950}
\draw[gray!60, dashed] (0,2.25) -- (0,1.0);
\draw[gray!60, dashed] ({(4950-1000)*\sx},2.25) -- (\W,1.0);
\foreach \i in {0,...,29} {
  \pgfplotstablegetelem{\i}{lo}\of\blocks \let\lo\pgfplotsretval
  \pgfplotstablegetelem{\i}{hi}\of\blocks \let\hi\pgfplotsretval
  \pgfplotstablegetelem{\i}{mainhi}\of\blocks \let\mh\pgfplotsretval
  \pgfplotstablegetelem{\i}{usedhi}\of\blocks \let\uh\pgfplotsretval
  \pgfmathsetmacro{\a}{(\lo-1000)*\zx}
  \pgfmathsetmacro{\u}{(\uh+1-1000)*\zx}
  \pgfmathsetmacro{\m}{(\mh+1-1000)*\zx}
  \pgfmathsetmacro{\b}{(\hi+1-1000)*\zx}
  \fill[flaggreen] (\a,0.4) rectangle (\u,1.0);
  \fill[flaggreen!25] (\u,0.4) rectangle (\m,1.0);
  \fill[pattern=north east lines, pattern color=amber] (\m,0.4) rectangle (\b,1.0);
  \draw[white, line width=0.6pt] (\a,0.4) -- (\a,1.0);
}
\draw[gray] (0,0.4) rectangle (\W,1.0);
% a few blocks named
\foreach \i/\side in {0/below, 8/above, 14/below, 20/above, 24/below, 29/above} {
  \pgfplotstablegetelem{\i}{lo}\of\blocks \let\lo\pgfplotsretval
  \pgfplotstablegetelem{\i}{hi}\of\blocks \let\hi\pgfplotsretval
  \pgfplotstablegetelem{\i}{zonename}\of\blocks \let\ar\pgfplotsretval
  \pgfmathsetmacro{\mid}{((\lo+\hi+1)/2-1000)*\zx}
  \ifx\side\above \draw[gray] (\mid,1.0) -- (\mid,1.25) node[above, font=\tiny, align=center, text=black] {\ar\\\lo–\hi};
  \else \draw[gray] (\mid,0.4) -- (\mid,0.15) node[below, font=\tiny, align=center, text=black] {\ar\\\lo–\hi}; \fi
}
% legend
\fill[flaggreen] (0,-0.95) rectangle (0.35,-0.7); \node[font=\footnotesize, anchor=west] at (0.4,-0.82) {\T{numbers issued}{numéros émis}};
\fill[flaggreen!25] (3.4,-0.95) rectangle (3.75,-0.7); \node[font=\footnotesize, anchor=west] at (3.8,-0.82) {\T{rest of the main part (infill)}{reste de la partie principale (insertion)}};
\fill[pattern=north east lines, pattern color=amber] (8.8,-0.95) rectangle (9.15,-0.7); \draw[gray] (8.8,-0.95) rectangle (9.15,-0.7);
\node[font=\footnotesize, anchor=west] at (9.2,-0.82) {\T{tail reserve (last third)}{réserve de queue (dernier tiers)}};
\end{tikzpicture}
\end{document}
